\documentclass[11pt]{article}
\usepackage[margin=1in]{geometry}
\usepackage{amsmath,amssymb,amsthm}
\usepackage{enumitem}
\usepackage{hyperref}

\title{Summary of Section~4: Convergence Analysis for QPCA-EnDCF}
\author{}
\date{}

\begin{document}
\maketitle

\section*{Purpose of the Section}

Section~4 establishes the main theoretical convergence guarantees for the
QPCA-EnDCF (Quantitative Principal Component Analysis--Ensemble Data-Consistent
Filtering) algorithm. Its central goal is to provide a rigorous bias--variance
decomposition for the mean-squared error (MSE) of the analysis ensemble mean,
making explicit how two distinct sources of error---spectral truncation and
finite-ensemble sampling---jointly determine performance.

The development proceeds in three stages:
\begin{enumerate}
  \item \textbf{Concentration estimates.} Baseline identities and moment-based
    bounds quantify how empirical covariance operators fluctuate around their
    population counterparts.
  \item \textbf{Spectral perturbation.} Deterministic matrix perturbation
    theory translates covariance deviations into stability results for
    eigenvalues and eigenspaces, controlling the accuracy of empirical spectral
    projectors.
  \item \textbf{Bias--variance decomposition.} The main theorem provides an
    explicit MSE decomposition that cleanly separates truncation effects
    (governed by the truncation level~$\kappa$ and alignment of the mean
    innovation with the retained subspace) from sampling effects (governed by
    ensemble size~$N$ and the stability of the empirical projector).
\end{enumerate}


%=======================================================================
\section*{Subsection 4.1: Regularity Assumptions}
%=======================================================================

\subsection*{Assumption~1 (Regularity conditions)}

This assumption collects four mild conditions that underpin the entire
convergence analysis:
\begin{enumerate}[label=(\roman*)]
  \item \textbf{Finite second moment.} The forecast distribution has
    $\mathbb{E}\|\mathbf{x}\|^2 < \infty$, ensuring all first- and
    second-moment quantities in the filter are well defined.
  \item \textbf{Positive-definite observation error covariance.} The stacked
    covariance $\mathbf{R}^{(L)}$ satisfies
    $\lambda_{\min}(\mathbf{R}^{(L)}) \ge r_{\min} > 0$, guaranteeing that
    the whitening transformation is well posed.
  \item \textbf{Finite fourth moment of whitened residuals.} The i.i.d.\
    whitened residuals satisfy
    $\mathbb{E}\|\mathbf{e}^{(1)}\|^4 < \infty$. This is strictly weaker
    than Gaussianity and provides the probabilistic input for all covariance
    concentration estimates.
  \item \textbf{Spectral decomposition.} The population covariance
    $\boldsymbol{\Sigma}_E$ admits a spectral decomposition with ordered
    eigenvalues. No spectral gap or decay condition is imposed unless
    explicitly stated later.
\end{enumerate}

\noindent
A remark on the scope of the i.i.d.\ assumption explains that while the formal
analysis requires independent ensemble members, practical cycling experiments
(where the analysis from one window becomes the forecast for the next) violate
this. Nevertheless, under ergodic dynamics and sufficient mixing, the
theoretical predictions---particularly the $\mathcal{O}(1/N)$ variance scaling
and the role of the cutoff gap $\delta_\kappa$---remain empirically valid.


%=======================================================================
\section*{Subsection 4.2: Fundamental Concentration Results}
%=======================================================================

This subsection establishes the elementary probabilistic identities and
concentration estimates that form the backbone of the convergence analysis.

\subsection*{Lemma~1 (Bias--variance decomposition for the analysis mean)}

This is the foundational identity: the MSE of the analysis mean decomposes
exactly into the squared deterministic bias and the variance:
\[
  \mathbb{E}\|\bar{\mathbf{x}}^a - \mathbf{x}^{\mathrm{true}}\|^2
  =
  \underbrace{\|\mathbb{E}[\bar{\mathbf{x}}^a] - \mathbf{x}^{\mathrm{true}}\|^2}
    _{\text{Bias}^2}
  +
  \underbrace{\mathbb{E}\|\bar{\mathbf{x}}^a - \mathbb{E}[\bar{\mathbf{x}}^a]\|^2}
    _{\text{Variance}}.
\]
The proof follows from the orthogonality of the deterministic bias vector and
the zero-mean fluctuation (i.e., the cross-term vanishes because
$\mathbb{E}[\boldsymbol{\xi}] = \mathbf{0}$). This identity is exact and holds
without any distributional assumptions beyond integrability.

\subsection*{Lemma~2 (Unbiasedness of the sample covariance)}

Under the i.i.d.\ sampling assumption and for $N \ge 2$,
$\mathbb{E}[\mathbf{C}_E] = \boldsymbol{\Sigma}_E$. That is, the Bessel-corrected
sample covariance is an unbiased estimator of the population covariance. This
ensures that deviations $\mathbf{C}_E - \boldsymbol{\Sigma}_E$ are purely
sampling fluctuations with zero mean.

\subsection*{Lemma~3 (Mean-square Frobenius deviation of the sample covariance)}

This lemma quantifies the rate at which the sample covariance concentrates
around the population covariance in the Frobenius norm:
\[
  \mathbb{E}\|\mathbf{C}_E - \boldsymbol{\Sigma}_E\|_F^2
  \le
  \frac{C_{\mathrm{cov}}}{N-1}\,
  \mathrm{tr}(\boldsymbol{\Sigma}_E^2).
\]
The bound scales as $\mathcal{O}(N^{-1})$ and is proportional to
$\mathrm{tr}(\boldsymbol{\Sigma}_E^2)$, which measures the ``effective
dimensionality'' of the covariance. The constant $C_{\mathrm{cov}}$ depends on
dimension~$d$ and standardized fourth moments.

\subsection*{Lemma~4 (Fourth-moment bound for the sample mean)}

For i.i.d.\ samples with finite fourth moment, the fourth moment of the
centered sample mean satisfies
\[
  \mathbb{E}\|\bar{\mathbf{e}} - \boldsymbol{\mu}_E\|^4
  \le
  \frac{C_{\mathrm{mean}}}{N^2}\bigl(\mathbb{E}\|\mathbf{e}^{(1)} -
    \boldsymbol{\mu}_E\|^4 + (\mathbb{E}\|\mathbf{e}^{(1)} -
    \boldsymbol{\mu}_E\|^2)^2\bigr).
\]
This gives a $(\mathbb{E}\|\cdot\|^4)^{1/2} = \mathcal{O}(N^{-1})$ rate. The
result is used later in Cauchy--Schwarz arguments to control the variance term
involving the product of projector error and sample-mean fluctuation.


%=======================================================================
\section*{Subsection 4.3: Spectral Perturbation Bounds}
%=======================================================================

This subsection converts matrix-level covariance deviations into spectral
stability estimates. The results are deterministic inequalities; randomness
enters only through the perturbation $\mathbf{C}_E - \boldsymbol{\Sigma}_E$.

\subsection*{Lemma~5 (Spectral perturbation and projector stability)}

This is a four-part lemma drawing on classical perturbation theory:
\begin{enumerate}[label=(\roman*)]
  \item \textbf{Weyl's inequality.} Each empirical eigenvalue
    $\hat{\lambda}_i$ satisfies
    $|\hat{\lambda}_i - \lambda_i| \le \|\mathbf{C}_E - \boldsymbol{\Sigma}_E\|_2$.
    That is, eigenvalue perturbations are bounded by the spectral norm of the
    covariance deviation.

  \item \textbf{Davis--Kahan $\sin\Theta$ theorem (individual eigenvectors).}
    If the $i$-th population eigenvalue $\lambda_i$ is isolated with spectral
    gap $\delta_i > 0$, then
    $\sin\angle(\hat{\mathbf{v}}_i, \mathbf{v}_i) \le
      \|\mathbf{C}_E - \boldsymbol{\Sigma}_E\|_2 / \delta_i$.
    Larger spectral gaps yield more stable eigenvector estimates.

  \item \textbf{Davis--Kahan for projectors.} If the cutoff gap satisfies
    $\lambda_\kappa - \lambda_{\kappa+1} \ge \delta_\kappa > 0$, then
    \[
      \|\hat{\mathbf{P}}_\kappa - \mathbf{P}_\kappa\|_F
      \le
      \frac{\sqrt{2\kappa}}{\delta_\kappa}\,
      \|\mathbf{C}_E - \boldsymbol{\Sigma}_E\|_F.
    \]
    This is the critical bound for QPCA-EnDCF, whose update is governed by the
    rank-$\kappa$ empirical projector.

  \item \textbf{Projected covariance perturbation.} The Frobenius deviation
    between the truncated empirical and population covariances decomposes via
    the triangle inequality into a ``direct'' covariance perturbation term and
    two ``cross'' terms involving $\hat{\mathbf{P}}_\kappa - \mathbf{P}_\kappa$.
\end{enumerate}

\noindent
A key insight is that spectral separation is intrinsic to eigenspace stability:
when eigenvalues are nearly repeated, small perturbations can cause large
rotations within the invariant subspace. For QPCA-EnDCF, projector stability is
strongest when the cutoff gap $\lambda_\kappa - \lambda_{\kappa+1}$ is large.

\subsection*{Lemma~6 (Fourth-moment bound for sample covariance error, Gaussian case)}

Under the additional assumption of Gaussian whitened residuals, the
fourth moment of the Frobenius covariance error satisfies
\[
  \mathbb{E}\|\mathbf{C}_E - \boldsymbol{\Sigma}_E\|_F^4
  \le
  \frac{C_{\mathrm{W}}}{(N-1)^2}\,
  \bigl(\mathrm{tr}(\boldsymbol{\Sigma}_E^2)\bigr)^2.
\]
The proof leverages the Wishart distribution of $(N-1)\mathbf{C}_E$ by showing
that the centering matrix is an orthogonal projector, so right-multiplication
by an orthogonal matrix preserves the Gaussian distribution of the data matrix.
This gives $(\mathbb{E}\|\cdot\|_F^4)^{1/2} = \mathcal{O}(N^{-1})$, which is
needed for sharp Cauchy--Schwarz arguments in the main theorem.


%=======================================================================
\section*{Subsection 4.4: Main Theorem---Bias--Variance Decomposition}
%=======================================================================

\subsection*{Theorem~1 (Bias--variance decomposition with spectral truncation)}

This is the central result of the section. Under the regularity assumptions and
a cutoff gap condition $\lambda_\kappa^{\,s} - \lambda_{\kappa+1}^{\,s} \ge
\delta_\kappa > 0$, the MSE decomposes exactly as
\[
  \mathbb{E}\|\bar{\mathbf{x}}^{a,s} - \mathbf{x}^{\mathrm{true}}\|^2
  =
  \mathrm{Bias}^2(\kappa, s) + \mathrm{Var}(N, \kappa, s).
\]

The theorem provides explicit upper bounds for both components:

\paragraph{Bias bound.} The squared bias is bounded by a sum of four terms:
\begin{enumerate}
  \item $\mathrm{Bias}_{\mathrm{base}}^2(s)$: the residual analysis bias after
    applying the (untruncated) population mean correction.
  \item A \textbf{truncation term} proportional to
    $\|(\mathbf{I} - \mathbf{P}_\kappa^{\,s})\boldsymbol{\mu}_E^{\,s}\|^2$,
    measuring how much of the mean innovation is discarded by restricting to
    rank~$\kappa$.
  \item A \textbf{map-approximation term} proportional to
    $\|Q_s - Q\|_{L^2(P_{\mathcal{X}})}^2$, capturing the error from using an
    approximate QoI map $Q_s$ instead of the exact $Q$.
  \item A \textbf{projector-estimation term} scaling as
    $\kappa / \delta_\kappa^2 \cdot \mathbb{E}\|\mathbf{C}_E^{\,s} -
      \boldsymbol{\Sigma}_E^{\,s}\|_F^2$,
    reflecting the cost of estimating the spectral projector from finite samples.
\end{enumerate}

\paragraph{Variance bound.} The variance is bounded by:
\begin{enumerate}
  \item An $\mathcal{O}(1/N)$ term from forecast-mean fluctuation.
  \item An $\mathcal{O}(1/N)$ term proportional to
    $\mathrm{tr}(\boldsymbol{\Sigma}_E^{\,s})$ from the sample-mean residual
    fluctuation.
  \item A \textbf{joint projector--sample-mean term} capturing the interaction
    between the projector estimation error and the sample-mean fluctuation.
    This term is controlled via Cauchy--Schwarz combined with the projector
    stability bound and fourth-moment estimates.
\end{enumerate}

\noindent
The pseudoinverse gain is essential in the small-ensemble regime ($N - 1 < d$),
and Assumption~2 (gain-moment condition) is the only probabilistic requirement
on the gain needed for these bounds.

A remark following the theorem notes that the truncation contribution depends on
the alignment of $\boldsymbol{\mu}_E^{\,s}$ with the discarded subspace and
that generic spectral tail bounds do not apply without additional structural
assumptions (e.g., source/RKHS-type conditions).


%=======================================================================
\section*{Subsection 4.5: Corollaries and Practical Implications}
%=======================================================================

\subsection*{Corollary~1 (Scaling of the projector-estimation term, Gaussian case)}

Under Gaussian whitened residuals, the joint projector--sample-mean variance
contribution scales as
\[
  \mathcal{O}\!\left(\frac{\kappa}{\delta_\kappa^2 \, N^2}\right).
\]
The proof combines the Cauchy--Schwarz-based marginal bound from the main
theorem with the Wishart fourth-moment bound (Lemma~6) and the sample-mean
fourth-moment bound (Lemma~4). The resulting $N^{-2}$ rate (rather than
$N^{-1}$) shows that this ``higher-order'' variance contribution is dominated
by the leading $\mathcal{O}(1/N)$ terms.

\subsection*{Corollary~2 (Comparison to stochastic EnKF observation perturbations)}

This corollary provides a direct theoretical comparison between QPCA-EnDCF and
the stochastic (perturbed-observation) EnKF. The stochastic update introduces
an irreducible variance floor:
\[
  \mathrm{Var}_{\mathrm{stoch}}
  \ge
  \frac{1}{N}\,
  \mathrm{tr}\!\bigl(\mathbf{K}^{(w)}\mathbf{R}^{(L)}(\mathbf{K}^{(w)})^\top\bigr),
\]
which scales as $\mathcal{O}(d/N)$ for isotropic observation noise
$\mathbf{R}^{(L)} = \bar{r}\,\mathbf{I}$. By contrast, QPCA-EnDCF avoids this
perturbation-induced variance source entirely. QPCA-EnDCF instead incurs
variability through estimation of the low-rank projector, but this is controlled
by the cutoff gap $\delta_\kappa$ and the covariance concentration rate, and
does not grow with observation-space dimension~$d$.

\subsection*{Corollary~3 (Principled selection of the truncation level)}

This corollary formulates the selection of the truncation level $\kappa$ as a
proxy optimization problem. The optimal $\kappa^\ast$ minimizes an upper-bound
objective that balances:
\begin{itemize}
  \item The \textbf{truncation remainder}
    $\mathcal{T}_\kappa^{\,s} = \|(\mathbf{I} -
      \mathbf{P}_\kappa^{\,s})\boldsymbol{\mu}_E^{\,s}\|^2$,
    which decreases as $\kappa$ increases (less information discarded).
  \item The \textbf{projector-estimation cost}
    $\kappa / \delta_\kappa^2$, which generally increases with $\kappa$ (more
    eigenvectors to estimate, potentially with smaller gaps).
\end{itemize}
The map-approximation and forecast-fluctuation terms are independent of
$\kappa$ and do not affect the minimizer. In practice, $\kappa^\ast$ is
determined by this truncation-vs-estimation tradeoff.


%=======================================================================
\section*{Summary of Main Contributions}
%=======================================================================

\begin{enumerate}
  \item \textbf{Exact bias--variance decomposition.} The MSE splits cleanly
    into squared bias and variance, with explicit dependence on the truncation
    level~$\kappa$, ensemble size~$N$, cutoff gap~$\delta_\kappa$, and
    QoI-map approximation quality.

  \item \textbf{$\mathcal{O}(1/N)$ variance scaling.} The canonical Monte
    Carlo rate is preserved, with a stability prefactor
    $\kappa / \delta_\kappa^2$ that quantifies the cost of empirical projector
    estimation.

  \item \textbf{Elimination of perturbation-induced variance.} Unlike the
    stochastic EnKF, QPCA-EnDCF does not inject observation noise into the
    analysis mean, removing an $\mathcal{O}(d/N)$ variance floor.

  \item \textbf{Principled truncation selection.} The theory provides a
    rigorous proxy objective for choosing $\kappa$ by balancing truncation bias
    against projector-estimation variance.

  \item \textbf{Mild assumptions.} The results require only finite fourth
    moments (weaker than Gaussianity) for the general bounds, with sharper
    rates available under Gaussian assumptions.
\end{enumerate}

\end{document}
